% Generated draft: requires mathematical and source review.
\documentclass[11pt]{article}
\usepackage[margin=25mm]{geometry}
\usepackage{fontspec}
\setmainfont{Latin Modern Roman}
\usepackage{amsmath,amssymb,mathtools,unicode-math}
\setmathfont{Latin Modern Math}
\usepackage{enumitem,needspace,xurl,hyperref}
\setlength{\parindent}{0pt}
\setlength{\parskip}{5pt}
\setlength{\emergencystretch}{3em}
\newcommand{\N}{\mathbb{N}}
\newcommand{\Z}{\mathbb{Z}}
\newcommand{\Q}{\mathbb{Q}}
\newcommand{\R}{\mathbb{R}}
\newcommand{\C}{\mathbb{C}}
\newcommand{\F}{\mathbb{F}}
\newcommand{\A}{\mathbb{A}}
\newcommand{\PP}{\mathbb{P}}
\newcommand{\E}{\mathbb{E}}
\newcommand{\eps}{\varepsilon}
\newcommand{\dd}{\,\mathrm{d}}
\begin{document}
\section*{UnsolvedMath ID: 30003222}
\textbf{Source problem code:} catalogue entry
\subsection{Definitions and mathematical statement}
The supplied question concerns a visible polar representation \((G,V)\). The record does not specify the ground field, the precise hypotheses on \(G\), or the convention for Hamiltonian reduction. The following customary complex reductive setting explains the notation conditionally; these conventions require confirmation against the original source.

Let \(G\) be a complex reductive algebraic group acting linearly on a finite-dimensional complex vector space \(V\), and put \(\mathfrak g=\operatorname{Lie}(G)\). Write
\[
V/\!/G=\operatorname{Spec}\C[V]^G,
\qquad \pi:V\longrightarrow V/\!/G.
\]
The representation is called \emph{visible} if every fiber of \(\pi\) contains only finitely many \(G\)-orbits. A vector is called semisimple if its \(G\)-orbit is closed. In the usual definition of a polar representation, there is a semisimple vector \(v\in V\) such that
\[
\mathfrak c_v=\{x\in V:\mathfrak g\cdot x\subseteq\mathfrak g\cdot v\},
\qquad
\dim\mathfrak c_v=\dim(V/\!/G).
\]
Such a subspace \(\mathfrak c=\mathfrak c_v\) is a Cartan subspace. Its associated finite group is
\[
W=N_G(\mathfrak c)/Z_G(\mathfrak c),
\]
where the numerator preserves \(\mathfrak c\) setwise and the denominator fixes it pointwise. Let \(\mathfrak c^*=\operatorname{Hom}_{\C}(\mathfrak c,\C)\), with the dual \(W\)-action. Interpret the left-hand quotient as the affine quotient
\[
(\mathfrak c\oplus\mathfrak c^*)/W
=\operatorname{Spec}\C[\mathfrak c\oplus\mathfrak c^*]^W,
\]
with its Poisson structure induced by the canonical symplectic structure on \(\mathfrak c\oplus\mathfrak c^*\).

Identify \(V\oplus V^*\) with \(T^*V\). Its cotangent moment map is
\[
\mu:V\oplus V^*\longrightarrow\mathfrak g^*,
\qquad
\mu(v,\lambda)(\xi)=\lambda(\xi\cdot v).
\]
If \(I_\mu\subseteq\C[V\oplus V^*]\) is the ideal generated by the functions \((v,\lambda)\mapsto\lambda(\xi\cdot v)\), for \(\xi\in\mathfrak g\), two possible zero-level reduction conventions are
\[
X_{\mathrm{sch}}=\operatorname{Spec}\bigl((\C[V\oplus V^*]/I_\mu)^G\bigr),
\qquad
X_{\mathrm{red}}=\operatorname{Spec}\bigl((\C[V\oplus V^*]/\sqrt{I_\mu})^G\bigr).
\]
The supplied notation \((V\oplus V^*)\mathbin{///}G\) does not specify which convention is intended. No equality of these two constructions is assumed here.

\textbf{Question (permanent catalogue entry).} For every visible polar representation \((G,V)\) in the intended source setting, is the natural Poisson morphism named in the record
\[
\phi:(\mathfrak c\oplus\mathfrak c^*)/W
\longrightarrow (V\oplus V^*)\mathbin{///}G
\]
an isomorphism of Poisson spaces? The record does not provide the construction of \(\phi\), including how the Cartan dual enters that construction. Its precise construction, the category of spaces, the reduction convention, and the ambient representation hypotheses remain qualifications requiring source review. Stability and local freeness are not assumptions in the supplied general question.
\subsection{Short English statement}
The question asks whether the natural Poisson map from the finite quotient of a Cartan subspace together with its dual to the symplectic reduction of \(V\oplus V^*\) is an isomorphism for every visible polar representation. The ground field, group hypotheses, reduction convention, and construction of the map require confirmation from the source.
\subsection{Sources}
\begin{itemize}
\item Supplied local UnsolvedMath registry record: permanent ID \(30003222\), title \emph{Symplectic Reduction of Visible Polar Representations}, collection \emph{Oberwolfach Reports Open Problems}, exact source problem code \texttt{catalogue entry}. The record gives proposed year \(2016\) and names no individual proposer.
\item Supplied report citation: \emph{Singularities (2017)}, Oberwolfach Reports, pp.\ 2625--2691, DOI \texttt{10.4171/owr/2016/46}. Original-conjecture link: \url{https://doi.org/10.4171/owr/2016/46}. The supplied background describes Oberwolfach Reports as reports of mathematical workshops and their open-problem discussions. The citation's \(2017\) label is retained alongside the DOI's \(2016\) component and the proposed year \(2016\).
\item Supplied literature source, labelled ``visible stable locally free case'': \url{https://arxiv.org/abs/1609.04573}. No authors, title, or theorem number are supplied for this reference.
\item Literature assessment dated exactly \texttt{2026-08-21}: the supplied curation labels the problem open and states that the proposed isomorphism is proved for visible stable locally free polar representations and several other families, while visibility alone does not yield a general theorem. This is supplied dated curation derived from the OWR record, not an independently verified literature assessment. Its status and cited claims require independent checking.
\item The registry assigns the statement status \texttt{corrected\_verified} and reports a manual semantic, typographic, and LaTeX review. These are supplied registry qualifications; they do not establish mathematical correctness or resolve the formulation issues identified here. No external source has been visited for this draft.
\end{itemize}
\end{document}
